% Generated draft: requires mathematical and source review.
\documentclass[11pt]{article}
\usepackage[margin=25mm]{geometry}
\usepackage{fontspec}
\setmainfont{Latin Modern Roman}
\usepackage{amsmath,amssymb,mathtools,unicode-math}
\setmathfont{Latin Modern Math}
\usepackage{enumitem,needspace,xurl,hyperref}
\setlength{\parindent}{0pt}
\setlength{\parskip}{5pt}
\setlength{\emergencystretch}{3em}
\newcommand{\N}{\mathbb{N}}
\newcommand{\Z}{\mathbb{Z}}
\newcommand{\Q}{\mathbb{Q}}
\newcommand{\R}{\mathbb{R}}
\newcommand{\C}{\mathbb{C}}
\newcommand{\F}{\mathbb{F}}
\newcommand{\A}{\mathbb{A}}
\newcommand{\PP}{\mathbb{P}}
\newcommand{\E}{\mathbb{E}}
\newcommand{\eps}{\varepsilon}
\newcommand{\dd}{\,\mathrm{d}}
\begin{document}
\section*{UnsolvedMath ID: 800167}
\textbf{Source problem code:} catalogue entry
\subsection{Definitions and mathematical statement}
For a Hecke--Maass cusp form \(f\), write \(L(s,f)\) for its standard \(L\)-function and \(L(s,\operatorname{ad}f)\) for its adjoint \(L\)-function. The spectral parameter \(t_f\) is related to the Laplace eigenvalue by \(\lambda_f=\tfrac14+t_f^2\). The source does not specify how exceptional imaginary spectral parameters are handled in its cutoff.

The supplied record attributes to Jutila the level-one twelfth-moment estimate
\[
\sum_{t_f\leq T}\frac{L\!\left(\tfrac12,f\right)^{12}}{L(1,\operatorname{ad}f)}
\ll_{\varepsilon}T^{4+\varepsilon},
\]
where the sum is over Hecke--Maass cusp forms of level \(1\), and the estimate is understood for every \(\varepsilon>0\) as \(T\to\infty\). The notation \(\ll_{\varepsilon}\) permits the implied constant to depend on \(\varepsilon\).

\textbf{Problem (Peter Humphries, Virginia, as attributed in the source).} Prove a level-aspect analogue of this Weyl-strength twelfth-moment estimate by summing over Hecke--Maass cusp forms of level \(q\), with spectral parameter bounded independently of the positive integer level \(q\), as \(q\to\infty\).

\textbf{Unresolved quantitative and family specifications.} The supplied question does not state the desired exponent of \(q\), restrictions on the levels, or the precise family at level \(q\). In particular, it does not specify the congruence subgroup, character, inclusion of oldforms, or multiplicity convention. It also does not explicitly prescribe the level-aspect weight or its normalization. Thus the source supplies a research question, but not a uniquely specified level-aspect inequality. These choices require clarification before a fully quantified numerical conjecture can be stated.
\subsection{Short English statement}
The question asks for an analogue, as the level grows and the spectral parameter stays bounded, of Jutila's stated level-one twelfth-moment bound. The desired power of the level and the precise family and weighting are unspecified. The registry labels the problem open; this draft does not independently verify that status.
\subsection{Sources}
\begin{itemize}
\item Local UnsolvedMath registry, supplied complete source record, permanent ID \(800167\), exact source code \texttt{catalogue entry}. The supplied record labels the problem \emph{open} and lists Hamieh, Kadiri, and Ng in its proposer field, with year \(2024\).
\item Hamieh, Kadiri, and Ng, \emph{Moments of L-functions Problem List} (2024), Problem 17, \url{https://arxiv.org/abs/2405.17800}. The supplied background records access on 2026-10-16 and extraction by \texttt{arxiv\_latex}; these are supplied provenance details, not evidence of independent source consultation for this draft.
\item The quoted question is attributed within the supplied statement to Peter Humphries, Virginia. The level-one estimate is attributed there to Jutila, but no separate bibliographic reference for that estimate is supplied.
\end{itemize}
\end{document}
